\chapter[Conclusiones]{Conclusiones}
\label{cp:conclusiones}
% Extensión estimada: 4--6 páginas. Sin información nueva.
% Feedback del tutor (14/09/2026): al eliminarse la sección "Contribuciones" de la introducción, los aportes
% del trabajo se enuncian acá, ya respaldados por los resultados.

\pendiente{Este capítulo cierra el trabajo con sus conclusiones y el grado de cumplimiento de los objetivos.}

\section{Conclusiones generales}
\label{sec:concl_generales}
\pendiente{Se sintetizan los hallazgos principales en pocos párrafos.}

\section{Cumplimiento de los objetivos}
\label{sec:concl_objetivos}
\pendiente{Una tabla relaciona cada objetivo específico con el capítulo que lo aborda y el resultado
obtenido.}

\section{Aportes del trabajo}
\label{sec:concl_aportes}
\pendiente{Se enuncian los aportes, cada uno con remisión al capítulo que lo sustenta: una comparación
sistemática bajo un mismo protocolo entre clasificadores supervisados, VLMs médicos generales y un modelo
visión--lenguaje dermatológico; el diagnóstico del colapso de los VLMs generales en clasificación
dermatológica y la evidencia de que el factor determinante es la especialización de dominio y no la
modalidad; un ensamble en cascada consciente de la incertidumbre que mejora la fidelidad diagnóstica de los
reportes, replicado con dos generadores; la auditoría de fuga de información entre particiones de HAM10000 y
los resultados corregidos con una partición disjunta por lesión; y el análisis de viabilidad sobre hardware
de consumo.}
